
\subsection*{4.4.}
On
appelle parfois l'entier (4.3.4) le \emph{nombre d'intersection} des @@M0@@ modules inversibles @@M1@@ avec @@M2@@, et on le dénote par @@M3@@ (cf. [3]). Comme conséquence de l'expression (4.3.4) de ce nombre on trouve aussitôt les propriétés suivantes :

\begin{enumerate}[label=\arabic*)]
\item L'entier (4.3.4) est une fonction symétrique et @@M4@@-linéaire dans les variables @@M5@@, nulle si @@M6@@, et il est «~additif~» en @@M7@@ i.e. additif en @@M8@@.

\item Utilisant 2.1.2, on en conclut immédiatement (avec les notations de 1.1.2) :
@@M9@@

\item Soit @@M10@@ un sous-schéma fermé de @@M11@@ et @@M12@@ un @@M13@@-module cohérent tel que @@M14@@, où @@M15@@ est l'immersion canonique. Alors on trouve
@@M16@@
De façon plus générale, la formule de projection donne en effet
@@M17@@
pour tout @@M18@@-morphisme @@M19@@ de @@M20@@-schémas propres.

\item Supposons que dans le morphisme @@M21@@ ci-dessus, @@M22@@ et @@M23@@ sont irréductibles de dimension @@M24@@. Si @@M25@@ est de degré @@M26@@, alors on a
